\section{Nekrasov residues for the $\mathbb C^3$ Rosetta and the toric CY$_3$ landscape}
\label{wn:c3-rosetta-nekrasov}

\providecommand{\Wilpha}{\mathcal{W}_{1+\infty}}
\providecommand{\gghone}{\widehat{\mathfrak{gl}}_1}
\providecommand{\CoHA}{\mathrm{CoHA}}
\providecommand{\Hilb}{\mathrm{Hilb}}
\providecommand{\Eone}{E_1}
\providecommand{\Etwo}{E_2}
\providecommand{\hCS}{\mathrm{hCS}}

The calculations take place over
$\mathbb F=\mathbb C(\epsilon_1,\epsilon_2)$ with
$\epsilon_3=-\epsilon_1-\epsilon_2$. The toric
$T^3$-equivariant input is scoped to geometries carrying such a torus;
compact non-toric CY$_3$ varieties have no $T^3$ action.

\subsection*{The $d=3$ functor for $\mathbb C^3$}

The $\mathbb C^3$ object splits into four statements.
\begin{enumerate}
\item $\CoHA(\mathbb C^3)\simeq Y^+(\gghone)$
    (Schiffmann--Vasserot 2013, Thm.\ 1.1 over
    $\mathbb F(\!(z)\!)$; positive half).
\item $Y(\gghone)\simeq\Wilpha[\lambda]$ at the self-dual level
    (Tsymbaliuk 2017, Thm.\ 1.1; full Drinfeld double).
\item $\Phi_3(\mathrm{Perf}(\mathbb C^3))$ is an $E_1$-chiral
    algebra.
\item The Drinfeld centre
    $\mathcal Z(\mathrm{Rep}^{\Eone}(-))$ carries the corresponding
    $\Etwo$-enhanced module category.
\end{enumerate}

The composite has the form
\[
\begin{array}{c}
\mathrm{Perf}(\mathbb C^3)
  \xrightarrow{\;\CoHA\;}
Y^+(\gghone)
  \xrightarrow{\;\mathcal D_\hbar\;}
Y(\gghone)
  \xrightarrow{\;\mathrm{Tsym}\;}
\Wilpha \\[4pt]
\hspace{3.3cm}\underbrace{\hspace{7.4cm}}_{=\;\Phi_3\text{ on this object}}
\end{array}
\]
and is verified one arrow at a time. The composite consists of three
separate identifications. The left-hand side is an $\Eone$-algebra over
$\mathbb F(\!(z)\!)$; the full $\Wilpha$ appears only after doubling,
Cartan data, completion, and vertex or defect evaluation. At $d\ge3$,
the specialised CY-to-chiral output is $\Eone$; the $\Etwo$ structure
lives on $\mathcal Z(\mathrm{Rep}^{\Eone})$.

\subsection*{$\CoHA(\mathbb C^3)$ and the positive half}

Schiffmann--Vasserot 2013 Thm.\ 1.1 identifies the critical CoHA of the
Jordan-triple-loop quiver with potential
$W=\mathrm{tr}(X[Y,Z])$ with the positive shuffle algebra:
\[
\CoHA(\mathbb C^3)\otimes_{\mathbb F}\mathbb F(\!(z)\!)
\simeq \mathrm{Sh}^+ \simeq Y^+(\gghone).
\]
The full affine Yangian is
\[
Y(\gghone)=Y^+(\gghone)\bowtie Y^0(\gghone)\bowtie Y^-(\gghone),
\]
the Hall--Drinfeld double with Cartan modes and a non-degenerate Hopf
pairing. Tsymbaliuk's current theorem identifies this completed double
with the affine Yangian presentation which acts on the corresponding
$\Wilpha[\lambda]$ vacuum representation. The raw CoHA supplies the
positive half.

The positive-half character is the MacMahon series
\[
\chi_{\CoHA(\mathbb C^3)}(q)
  = \sum_{n\ge0}q^n\,\chi_T(\Hilb^n(\mathbb C^2))
  = M(q)
  = \prod_{n\ge1}\frac{1}{(1-q^n)^n}.
\]
At the self-dual level $c=1$, the Heisenberg specialisation of the full
$\Wilpha$ has character
\[
\chi_{\Wilpha,c=1}(q)
  = \chi_{H_1}(q)
  = P(q)
  = \prod_{n\ge1}\frac{1}{1-q^n}.
\]
The quotient
\[
\frac{M(q)}{P(q)}=\prod_{n\ge2}(1-q^n)^{-(n-1)}
\]
records the positive-half tower of higher-spin modes. This is compatible
with the bar-Euler computation
$\sum_k(-1)^k\mathrm{ch}(B^k)=1/M(q)$
(Proposition~\ref{wn:prop:c3-bar-euler}).

\begin{center}
\renewcommand{\arraystretch}{1.15}
\begin{tabular}{lll}
\toprule
$Y^+(\gghone)$ generator & shuffle representative &
equivariant-cohomology class \\
\midrule
$e_k$ & $p_k(z)=\sum_i z_i^k$
  & equivariant $k$-th power sum on $\Hilb^\bullet(\mathbb C^2)^T$ \\
$\psi_k$ & $\omega^{\star k}$ vacuum-to-vacuum mode
  & Chern character of the tautological bundle \\
$f_k$ & absent from $Y^+$; appears in $Y^-$ after $\mathcal D_\hbar$
  & Serre-dual generator in the double \\
\bottomrule
\end{tabular}
\end{center}

The positive-half character has three independent witnesses:
\begin{enumerate}
\item localisation on $\Hilb^n(\mathbb C^2)^T$ with the CY$_3$ weight
    lift $\epsilon_3=-\epsilon_1-\epsilon_2$;
\item MacMahon's plane-partition product
    $\prod_{n\ge1}(1-q^n)^{-n}$;
\item the degree-zero DT partition function
    $Z_{DT}(\mathbb C^3)=M(-q)^{\chi(\mathbb C^3)}$
    (MNOP I 2006; topological vertex of Aganagic--Klemm--Mariño--Vafa
    2005).
\end{enumerate}

\subsection*{The Drinfeld centre and the $\Etwo$ enhancement}

The centre statement is
\[
\mathcal Z\bigl(\mathrm{Rep}^{\Eone}(Y^+(\gghone))\bigr)
\simeq
\mathrm{Rep}^{\Etwo}(Y(\gghone))
\simeq
\mathrm{Rep}^{\Etwo}(\Wilpha).
\]
Müger's theorem gives a braided tensor category, hence an $\Etwo$
object, from the Drinfeld centre of a finite tensor category; the
pro-finite Yangian version is the corresponding completed centre. The
chiral algebra itself remains $\Eone$ on a reference curve. The
braiding belongs to the module category.

The centre character is
\[
\mathrm{ch}\,\mathcal Z(\mathrm{Rep}(Y^+))
  =M(q)^2P(q)
  =\prod_{n\ge1}(1-q^n)^{-(2n+1)}.
\]
Here the two factors of $M(q)$ are left and right positive-half module
data, and $P(q)$ is the neutral Cartan partition data. Direct
enumeration of half-braidings on Fock modules through rank $n\le10$
and Feigin--Tsymbaliuk 2011 (\texttt{arXiv:1101.0055}) give the same
elliptic-Hall lift. On the charge-$2$ sector, the Yang matrix
\[
R(u)=1+\hbar P/u+O(\hbar^2)
\]
satisfies Yang--Baxter and unitarity
$R(u)R(-u)=1$, with the CY$_3$ input
$\prod_i(z-\epsilon_i)(z+\epsilon_i)=\prod_i(z^2-\epsilon_i^2)$.
This is a centre calculation.

\subsection*{The shadow coefficient $m_5=-80/9$}

The shadow coefficients used in this note are
\[
(m_5,m_7,m_8)=\left(-\frac{80}{9},-\frac{24320}{81},
\frac{33157760}{19683}\right).
\]
The independent connected Wick value is
\[
G_5^{\mathrm{conn}}(1,2,3,4,5)=\frac{775}{5184}.
\]
The constants $S_5$ and $G_5^{\mathrm{conn}}$ have different
normalisations. The bridge from the connected five-point
function to $m_5$ passes through the Virasoro normalisation
$\langle T,T\rangle=c/2$, the cubic coefficient $\alpha=2$, and the
Zamolodchikov level-$2$ constant $S_4=10/27$ at $c=1$.

The Newton recursion gives
\[
a_3
=-\frac{a_1a_2+a_2a_1}{2a_0}
=-\frac{2\cdot6\cdot(40/27)}{2}
=-\frac{80}{9},
\]
with $(a_0,a_1,a_2)=(1,6,40/27)$. Thus
\[
m_5(T,T,T,T,T)=a_3T=-\frac{80}{9}T,
\qquad
S_5=\frac{a_3}{5}=-\frac{16}{9}.
\]
The raw product $-120\cdot775/5184=-3875/216$ is a distinct Wick
normalisation. A graph proof of $m_5=-80/9$ must exhibit the
normalisation map from the Wick correlator to the Newton coefficient
$a_3$.

The graph enumeration itself is
\[
G_5^{\mathrm{conn}}
=\sum_{[G]:\,|E(G)|-|V(G)|+1=5}
\frac{w(G)}{|\mathrm{Aut}(G)|}
=\frac{775}{5184}.
\]
The five connected trivalent graph classes carry symmetry factors
\[
\left\{\frac1{16},\frac18,\frac1{24},\frac1{48},\frac1{384}\right\}.
\]
The rational MZV projection has
$\zeta^{\mathrm{sv}}(2)=0$, eliminating spurious $\zeta(2)$ terms and
forcing the coefficient into $\mathbb Q$.

\subsection*{Conifold and local $\mathbb P^2$}

The conifold is local CY$_3$ and toric, with compact curve
$\mathbb P^1$. Its partition function is geometry-specific:
\[
Z_{DT}(\mathrm{conifold},q,Q)=
M(-q)^2\prod_{n\ge1}(1-(-q)^nQ)^n(1-(-q)^nQ^{-1})^n
\]
(Szendrői 2008; Morrison--Mozgovoy--Nagao--Szendrői 2010). The
wall-crossing between the two DT chambers is the Kontsevich--Soibelman
pentagon in $U_q(\mathfrak{sl}_2)$.

Let $Q_{\mathrm{con}}$ be the conifold quiver and
$J_{\mathrm{con}}=\mathbb CQ_{\mathrm{con}}/\partial W$ its Jacobi
algebra, with
$W=\mathrm{tr}(a_1b_1a_2b_2-a_1b_2a_2b_1)$. The McKay correspondence at
dimension vector $(1,1)$ gives two stable simples
$S_\circ,S_\bullet$. Their equivariant $K$-theory classes generate a
rank-$2$ Heisenberg algebra, and the compact $\mathbb P^1$ contribution
leaves
\[
\kappa_{\mathrm{ch}}(\mathrm{conifold})=1.
\]
This value comes from the direct McKay calculation.

For $X=\mathrm{Tot}(\mathcal O_{\mathbb P^2}(-3))$, the Beilinson
quiver has three vertices, nine arrows, and the CY$_3$ fibre potential.
At dimension vector $(1,1,1)$,
\[
\kappa_{\mathrm{ch}}(\mathrm{local}\,\mathbb P^2)
  =\frac{\chi(\mathbb P^2)}{2}
  =\frac32.
\]
This agrees with Gopakumar--Vafa growth
$|n^0_d|\sim27^d$, equivalently radius $R=1/27$ and shadow class
$\mathbf M$.

\subsection*{$\Omega$-background scope}

Nekrasov's $\Omega$-background is defined on toric geometries admitting
a torus action with fixed-point localisation.
\begin{itemize}
\item $\mathbb C^3$: $T^3$-toric.
\item Conifold: $T^3$-toric perturbatively; compact-curve
    wall-crossing is additional data.
\item Local $\mathbb P^2$: $T^3$-toric.
\item Local $\mathbb P^1\times\mathbb P^1$: $T^3$-toric.
\end{itemize}
The geometries $K3\times E$, the quintic, and compact CY$_3$ varieties
in general require separate compact methods. Transferring
$\Omega$-background localisation, 3d-partition sums, Nekrasov partition
functions, or MacMahon products to compact non-toric CY$_3$ varieties
requires an explicit torus action and a separate localisation theorem.

\subsection*{The $\mathbb C^3$ Rosetta Stone}

\begin{theorem}[$\mathbb C^3$ Rosetta Stone]
\label{wn:thm:c3-rosetta-nekrasov}
Let $X=\mathbb C^3$ with $T^3$-action of weights
$(\epsilon_1,\epsilon_2,\epsilon_3)$ satisfying
$\epsilon_1+\epsilon_2+\epsilon_3=0$. Then:
\begin{enumerate}[label=\textup{(\alph*)}]
\item $\CoHA(X)\simeq Y^+(\gghone)$ as associative algebras over
    $\mathbb F(\!(z)\!)$ (Schiffmann--Vasserot 2013).
\item The Hall--Drinfeld double
    $\mathcal D_\hbar(Y^+)=Y(\gghone)$ is the completed affine Yangian
    identified by Tsymbaliuk 2017.
\item After the chosen defect or vacuum evaluation, the $d=3$ chiral
    object is the $\Eone$ Heisenberg VOA $H_1$.
\item The Drinfeld centre
    $\mathcal Z(\mathrm{Rep}^{\Eone}(Y^+(\gghone)))$ is $\Etwo$, has
    character $M(q)^2P(q)$, and carries the Yang matrix
    $R(u)=1+\hbar P/u+O(\hbar^2)$.
\item $\kappa_{\mathrm{ch}}(X)=\kappa_{\mathrm{ch}}(H_1)=1$, witnessed
    by the OPE $J(z)J(w)\sim1/(z-w)^2$, genus-$1$ free energy
    $F_1=1/24$, DT degree-zero $F_1=1/24$, the affine-Yangian structure
    function at the self-dual point, and the CY categorical trace.
\item $m_5=-80/9$ is the Newton coefficient
    $a_3=-(a_1a_2+a_2a_1)/(2a_0)$ with
    $(a_0,a_1,a_2)=(1,6,40/27)$; the independent Wick witness is
    $G_5^{\mathrm{conn}}(1,2,3,4,5)=775/5184$.
\end{enumerate}
\end{theorem}

\begin{remark}[Proof obligations for full equivalence]
The full equivalence with $\Wilpha$ requires the Drinfeld double,
Cartan modes, a non-degenerate Hopf pairing, a specified completion,
and a vertex or defect evaluation. None of these data is present in
the raw CoHA. At $d\ge3$, an $\Etwo$ enhancement belongs to the centre
or doubled module category. The specialised output is $\Eone$.
\end{remark}
